\section{خواندن اولین دما با سنسور}
\subsection{قرار است چه چیزی بسازیم؟}
کار اول ساده است: یک سنسور را وصل کنیم و دمای آن را در کیل ببینیم. وقتی این مرحله درست شد، همان تابع خواندن را در پروژهٔ دیگر استفاده می‌کنیم. بعد سراغ داخل درایور می‌رویم. برای انجام مراحل اول لازم نیست اشاره‌گر تابع یا پروتکل باس را بلد باشید.

این نسخه برای
\textbf{یک سنسور با تغذیهٔ مستقل در هر باس}
است. دما را به سلسیوس می‌خواند و خطا را بررسی می‌کند. رزولوشن تغییر نمی‌کند؛ جست‌وجوی چند سنسور، نوشتن EEPROM و تغذیهٔ انگلی پیاده‌سازی نشده‌اند تا نسخهٔ اول قابل فهم بماند.

\checkpoint{\textbf{نتیجهٔ این فصل:} مدار و نمونه را اجرا می‌کنید، نمایش ثابت Watch را از خطای سنسور تشخیص می‌دهید، از تابع خواندن استفاده می‌کنید، پایه را در CubeMX تغییر می‌دهید و می‌فهمید بایت‌های دما چگونه به یک عدد تبدیل می‌شوند.}

مرجع سخت‌افزار،
\extlink{https://www.analog.com/media/en/technical-documentation/data-sheets/ds18b20.pdf}{دیتاشیت رسمی سنسور}
با شمارهٔ
\lr{19-7487, Rev. 6, July 2019}
است
\cite{ds18b20}.

\subsection{مرحلهٔ ۱: سه سیم و یک مقاومت را بررسی کنیم}
نمونهٔ موجود از
\lr{STM32F107VCT6}
استفاده می‌کند. سنسور را طبق جدول ببندید. پیش از انتخاب پایهٔ فیزیکی VDD و GND، شکل پکیج را ببینید؛ نمای پایین قطعه را نباید نمای بالا در نظر گرفت.
\begin{center}
\begin{tabular}{p{38mm}p{91mm}}
\toprule
سیگنال سنسور & اتصال به برد\\
\midrule
\lr{VDD} & تغذیهٔ ۳٫۳ ولت\\
\lr{GND} & همان زمین میکروکنترلر\\
\lr{DQ} & پایهٔ PB10 برای داده\\
مقاومت pull-up & ۴٫۷ کیلو اهم بین DQ و ۳٫۳ ولت\\
\bottomrule
\end{tabular}
\end{center}
DQ تنها سیم داده است. GPIO باید open drain باشد؛ یعنی میکروکنترلر یا خط را پایین بکشد یا آن را رها کند. مقاومت خط رهاشده را بالا می‌برد. در این نمونه، pull-up داخلی GPIO جای مقاومت بیرونی نیست.

کلاک پروژه، کریستال خارجی ۸ مگاهرتز و کلاک پردازندهٔ ۷۲ مگاهرتز است. با کریستال واقعی برد تطبیق دهید. اتصال SWD از PA13 و PA14، زمین و ولتاژ مرجع پروگرامر تشکیل می‌شود. برای تست اول، نمایشگر یا ترمینال سریال لازم نداریم.

\checkpoint{\textbf{پیش از ادامه:} یک سنسور داریم، VDD مستقل وصل است، زمین مشترک است، DQ به PB10 می‌رود و مقاومت بیرونی وجود دارد.}

\subsection{مرحلهٔ ۲: بازکردن، ساخت و دانلود در کیل}
\begin{enumerate}
\item فایل
\lr{ds18b20/MDK-ARM/DS18B20.uvprojx}
را باز کنید. اگر پروژه را جای دیگری کپی کرده‌اید، فایل همان کپی را باز کنید.
\item پک STM32F1 نسخهٔ 2.4.1 و Compiler 6 نصب باشند. کامپایلر بررسی‌شده در این کار نسخهٔ 6.24 است.
\item در
\lr{Options for Target}
و بخش Debug پروگرامر واقعی را انتخاب کنید. برای J-Link همان J-Link و برای ST-Link درایور مناسب آن را انتخاب کنید. برای سنسور واقعی، Simulator انتخاب نشود.
\item در Utilities بررسی کنید پروگرام با همان پروگرامر انجام شود. Build بگیرید و خطاهای ساخت را قبل از دیباگ برطرف کنید.
\item firmware جدید را Download کنید و وارد Debug شوید. هر بار C عوض شد، دوباره Build و Download لازم است؛ ویرایش فایل، برنامهٔ روی برد را عوض نمی‌کند.
\item breakpointهای مسیر خواندن را حذف یا غیرفعال و برنامه را آزاد اجرا کنید.
\end{enumerate}
برنامه یک بار راه‌اندازی و بارها خواندن را انجام می‌دهد. تبدیل سنسور ممکن است حدود ۷۵۰ میلی‌ثانیه طول بکشد. صفر اولیه فقط مقدار شروع متغیر C است و اندازه‌گیری موفق نیست.

برای PlatformIO پوشهٔ پروژه را باز و محیط پروگرامر مناسب را اجرا کنید:
\begin{LTR}
\begin{lstlisting}[language={},numbers=none]
pio run -e stm32f107vc_jlink
pio run -e stm32f107vc_jlink -t upload
pio debug -e stm32f107vc_jlink
# For an ST-Link probe, use stm32f107vc_stlink.
\end{lstlisting}
\end{LTR}
PlatformIO Core باید نصب باشد. ساخت اول toolchain مشخص‌شده را دریافت می‌کند. firmware از HAL، CMSIS، startup و linker محلی همین پروژه استفاده می‌کند.

\subsection{مرحلهٔ ۳: Watch برنامهٔ در حال اجرا را نشان بدهد}
در جلسۀ Debug، Watch 1 را باز کنید و با دوبار کلیک روی ردیف خالی نام‌های زیر را اضافه کنید. از
\textbf{نام متغیر سراسری}
\lr{g\_temperature\_c}
استفاده کنید، نه متغیر محلی
\lr{value}
داخل تابع.
\begin{center}
\begin{tabular}{p{57mm}p{72mm}}
\toprule
عبارت & چه چیزی نشان می‌دهد؟\\
\midrule
\lr{\texttt{g\_temperature\_c}} & آخرین دمای موفق\\
\lr{\texttt{g\_temperature\_status}} & نتیجهٔ آخرین تلاش کامل\\
\lr{\texttt{g\_temperature\_valid}} & موفق‌بودن همان تلاش\\
\lr{\texttt{g\_measurement\_busy}} & در حال انجام بودن خواندن جدید\\
\lr{\texttt{g\_successful\_samples}} & تعداد خواندن‌های واقعاً موفق\\
\lr{\texttt{g\_sample\_cycles}} & تعداد تلاش‌های کامل، حتی ناموفق\\
\lr{\texttt{g\_last\_success\_ms}} & تیک HAL در آخرین موفقیت\\
\bottomrule
\end{tabular}
\end{center}
گزینهٔ
\lr{View $\rightarrow$ Periodic Window Update}
را فعال کنید. کیل هنگام توقف اجرای برنامه، Watch را تازه می‌کند؛ به‌روزرسانی حین اجرا به این گزینه نیاز دارد
\cite{keilwatch}.
اگر خاموش باشد، breakpoint باعث دیده‌شدن مقدار جدید می‌شود چون توقف، نمایش را تازه می‌کند. این مشاهده ثابت نمی‌کند که breakpoint باعث اندازه‌گیری شده است.

هنگام بررسی این مشکل، همین گزینه در فایل تنظیمات محلی صفر بود. به یک تغییر داده شد و نام‌های Watch اضافه شدند. جلسۀ باز کیل، تنظیماتی را که قبلاً بارگذاری کرده نگه می‌دارد؛ در همان جلسۀ باز هم گزینهٔ منو را فعال کنید. فایل محلی
\lr{.uvoptx}
وارد گیت نمی‌شود. برای تکرار تنظیمات، پروژه را ذخیره و کیل را ببندید، سپس از پوشهٔ پروژه اجرا کنید:
\begin{LTR}
\begin{lstlisting}[language={},numbers=none]
python tools/setup_keil_watch.py
\end{lstlisting}
\end{LTR}
این ابزار تنظیمات را پشتیبان می‌گیرد و انتخاب پروگرامر، Watchهای قبلی و breakpointها را حفظ می‌کند. اگر جلسۀ باز بعداً تنظیمات قدیمی را ذخیره کرد، پس از بستن کیل ابزار را دوباره اجرا کنید.

برای مثال، دمای ۲۴٫۵، وضعیت صفر، valid برابر true و شمارندهٔ موفق رو به افزایش، یک نتیجۀ عادی موفق است. عدد دمای واقعی به مدار و محیط وابسته است. در اتاق ثابت ممکن است دما ثابت بماند؛
\textbf{شمارندهٔ موفق}
نشان می‌دهد خواندن‌های جدید کامل می‌شوند یا نه.

اعتبار به آخرین تلاش
\emph{کامل‌شده}
مربوط است و هنگام تبدیل بعدی دیده می‌ماند. busy همان تبدیل جدید را جدا نشان می‌دهد. اگر تلاش کامل‌شده خطا بدهد، valid برابر false و عدد قدیمی در حافظه باقی می‌ماند. در خطا از آن عدد به‌عنوان اندازه‌گیری تازه استفاده نکنید.

\checkpoint{\textbf{معیار عبور از این مرحله:} در اجرای آزاد، شمارندهٔ موفق زیاد شود و وضعیت صفر بماند. اگر فقط صفحه ثابت است، پیش از تغییر درایور، refresh و نام متغیر سراسری را بررسی کنید.}

\subsection{مرحلهٔ ۴: نمایش خراب است یا خواندن سنسور؟}
\begin{enumerate}
\item مقدار فقط با breakpoint ظاهر می‌شود: اول تنظیم به‌روزرسانی دوره‌ای را بررسی کنید. صفحهٔ ثابت و CPU متوقف یک چیز نیستند.
\item شمارندهٔ موفق زیاد می‌شود ولی دما تقریباً ثابت است: سنسور را کمی گرم کنید و چند نمونه را مقایسه کنید.
\item تلاش‌ها زیاد می‌شوند ولی موفقیت‌ها نه: کد وضعیت را ببینید. حلقه اجرا می‌شود ولی تراکنش سنسور خطا دارد.
\item دسترسی زنده هنوز ممکن نیست: برنامه را اجرا کنید و در فاصلهٔ ۲۵۰ میلی‌ثانیه‌ای بین تلاش‌ها Pause کنید. دو مشاهدهٔ شمارنده را مقایسه کنید. بعضی تنظیمات پروگرامر یا جلسۀ دیباگ ممکن است دسترسی یکسان به حافظهٔ حین اجرا ندهند.
\item شمارنده در توقف‌های امن هم پیش نمی‌رود: محل program counter را ببینید. ماندن در Error Handler یعنی مشکل راه‌اندازی؛ fault handler یعنی استثنای CPU. ماندن در HAL Delay به SysTick فعال و وقفهٔ باز نیاز دارد. فقط از عدد Watch نتیجه نگیرید درایور گیر کرده است.
\end{enumerate}
برای دیدن دمای زنده، داخل تایم‌اسلات GPIO، تأخیر میکروثانیه یا polling توقف نگذارید. مکث دیباگر، ارتباط زمان‌دار را به‌هم می‌زند. breakpoint پس از خواندن کامل برای یک مشاهده مفید است، ولی روش معمول نمایش نیست.

\begin{center}
\begin{tabular}{p{35mm}p{94mm}}
\toprule
وضعیت & معنی و اولین اقدام\\
\midrule
صفر، موفق & دما قابل استفاده است.\\
۱، آرگومان نامعتبر & اشاره‌گر خروجی و اتصال callbackها بررسی شود.\\
۲، نبود قطعه & VDD، زمین، DQ، مقاومت و پایهٔ واقعی بررسی شود.\\
۳، timeout & تغذیهٔ مستقل، کلاک و تأخیر بررسی شود.\\
۴، خطای CRC & اتصال، نویز، rise time، زمان‌بندی و توقف دیباگر بررسی شود.\\
۵، دادهٔ نامعتبر & نوع قطعه، بایت‌های ثابت و خط پایین بررسی شود.\\
\bottomrule
\end{tabular}
\end{center}

\subsection{مرحلهٔ ۵: دو تابع در پروژهٔ STM32 خودمان}
آداپتر سخت‌افزار، GPIO و تأخیر را پشت دو تابع نگه می‌دارد. تابع
\lr{DS18B20\_STM32\_Init()}
یک بار آماده‌سازی می‌کند و تابع
\lr{DS18B20\_STM32\_ReadTemperature(\&temperature)}
هر بار یک تبدیل جدید را شروع و نتیجه را می‌خواند.
تابع
\textbf{دو خروجی}
دارد: مقدار برگشتی می‌گوید عملیات موفق بوده یا نه؛ خود دما با اشاره‌گر نوشته می‌شود. عبارت
\lr{\&temperature}
یعنی «آدرس این متغیر». نوشتن
\lr{temperature = DS18B20\_STM32\_ReadTemperature(...)}
اشتباه است چون کد وضعیت، نه دما، در متغیر قرار می‌گیرد.

نوع
\lr{float}
عدد اعشاری را نگه می‌دارد. نوع
\lr{DS18B20\_Status}
نتیجۀ عملیات و ثابت
\lr{DS18B20\_OK}
موفقیت را نشان می‌دهد. بخش if فقط بعد از موفقیت از دما استفاده می‌کند.

\begin{LTR}
\begin{lstlisting}
#include "ds18b20_stm32.h"

/* After HAL, clock and GPIO initialization: */
DS18B20_STM32_Init();

/* From a foreground loop, with interrupts enabled: */
float temperature;
DS18B20_Status result = DS18B20_STM32_ReadTemperature(&temperature);
if (result == DS18B20_OK) {
    /* temperature can now be used. */
} else {
    /* Report the failure; do not use temperature here. */
}
\end{lstlisting}
\end{LTR}
این مثال کوتاه جای فراخوانی‌ها را نشان می‌دهد: include در بخش هدر، راه‌اندازی بعد از تنظیم برد و خواندن در حلقه. این قطعه، یک main کامل برای جایگزین‌کردن فایل تولیدشده نیست.



برای اضافه‌کردن به پروژهٔ STM32F107 دیگر:
\begin{enumerate}
\item دو پوشهٔ قابل‌حمل
\lr{Libraries/OneWire}
و
\lr{Libraries/DS18B20}
را کپی کنید.
\item فایل‌های
\lr{Application/Inc/ds18b20\_stm32.h}
و
\lr{Application/Src/ds18b20\_stm32.c}
را کپی کنید.
\item دو فایل C کتابخانه و فایل C آداپتر را به ساخت اضافه کنید. سه پوشهٔ Inc مربوط به آن‌ها در مسیر هدر کامپایلر باشد.
\item برچسب GPIO خروجی open drain را در CubeMX برابر
\lr{ONE\_WIRE\_1}
قرار دهید. آداپتر ماکروهای تولیدشدهٔ پایه را می‌خواند. کلاک درست باشد و هنگام خواندن وقفه فعال بماند.
\item پس از تنظیم برد، یک بار Init و در حلقهٔ معمول برنامه ReadTemperature را با بررسی وضعیت اجرا کنید.
\end{enumerate}
این آداپتر برای همین Cortex-M3 و رجیسترهاست. MCU دیگر آداپتر، vendor، startup و linker مناسب خودش را می‌خواهد؛ دو کتابخانهٔ قابل‌حمل به هدر STM32 نیاز ندارند.

\subsection{مرحلهٔ ۶: نمونهٔ کامل اجرایی را بفهمیم}
کد زیر همان برنامه‌ای است که firmware نمونه اجرا می‌کند:
\begin{LTR}
\lstinputlisting{../../ds18b20/Application/Src/app.c}
\end{LTR}
از این ترتیب بخوانید: تابع راه‌اندازی، Init آداپتر را صدا می‌زند. پردازش، busy را فعال می‌کند، دما می‌خواهد، وضعیت را ثبت و فقط نتیجۀ موفق را منتشر می‌کند. سپس شمارندهٔ تلاش کامل را زیاد و busy را پاک می‌کند و پیش از تلاش بعدی مکث دارد.

متغیر موقت
\lr{value}
فقط برای همان فراخوانی است. متغیرهای
\lr{g\_...}
سراسری‌اند و Watch آدرس ثابت آن‌ها در RAM را پیدا می‌کند. volatile تضمین می‌کند کامپایلر نوشتن را انجام دهد؛ گزینهٔ refresh نرم‌افزار را فعال نمی‌کند. این چند فیلد جدا نیز یک پیام همگام RTOS نیستند.

در موفقیت، شمارندهٔ موفق و تیک آخرین موفقیت عوض می‌شوند. در خطا، فقط شمارندهٔ تلاش و وضعیت و اعتبار عوض می‌شوند؛ دمای قبلی و زمان آخرین موفقیت دست‌نخورده می‌مانند. تست جدید برنامه همین توالی و برگشت از خطا را بررسی می‌کند.

main تولیدشده فقط سه اتصال به برنامه دارد:
\begin{LTR}
\begin{lstlisting}
/* USER CODE BEGIN Includes */
#include "app.h"
/* USER CODE END Includes */

/* USER CODE BEGIN 2 */
App_Init();
/* USER CODE END 2 */

/* USER CODE BEGIN 3 */
App_Process();
/* USER CODE END 3 */
\end{lstlisting}
\end{LTR}
اولی در بخش include، دومی پس از تنظیم HAL و کلاک و GPIO و سومی داخل حلقهٔ بی‌نهایت تولیدشده است. آکولادهای حلقه حفظ شوند. یا نمونهٔ پردازش موجود را استفاده کنید یا حلقهٔ خواندن خودتان را؛ دو خواننده هم‌زمان روی یک باس اجرا نکنید.

\subsection{مرحلهٔ ۷: تغییر پایه یا شروع از کپی پروژه}
برای نمونهٔ مستقل کامل، کل پوشهٔ پروژهٔ سنسور را کپی کنید. برای درایور جدید، قالب خام را کپی کنید؛ از آن به‌عنوان وابستگی ساخت استفاده نکنید. هر پروژه IOC، vendor و تنظیمات ساخت محلی دارد.

برای تغییر GPIO همین نمونه:
\begin{enumerate}
\item IOC را باز کنید. پایه را جابه‌جا کنید ولی برچسب
\lr{ONE\_WIRE\_1}
و خروجی open drain، مقدار اولیهٔ high، بدون pull و سرعت high را حفظ کنید.
\item کلاک و SWD درست بمانند. در Project Manager حفظ user code، کپی vendor محلی، MDK-ARM V5 و نام پروژهٔ DS18B20 تنظیم شوند. Generate در همان پوشه انجام شود
\cite{cubemx}.
\item سه اتصال برنامه در USER CODE را کنترل کنید. فایل‌های Application و Libraries از Core تولیدشده جدا هستند.
\item دستورهای زیر، سپس Build کیل و Download firmware جدید را انجام دهید.
\end{enumerate}
\begin{LTR}
\begin{lstlisting}[language={},numbers=none]
python tools/sync_keil.py
python tools/check_project.py
python tests/run_tests.py
pio run -e stm32f107vc_jlink -e stm32f107vc_stlink
\end{lstlisting}
\end{LTR}
hook بعد از تولید تلاش می‌کند sync و check را خودکار انجام دهد. اگر Python در محیط CubeMX نیست، دستی اجرا کنید. sync، فایل‌های برنامه، آداپتر و کتابخانه را به کیل اضافه می‌کند؛ فایل پروژهٔ تولیدشده به‌تنهایی تمام فایل‌های C دست‌نویس را نمی‌شناسد.

\checkpoint{\textbf{تست انتقال پایهٔ انجام‌شده:} در کپی موقت، PB10 به PB11 منتقل و پروژه Generate و با هر دو toolchain ساخته شد؛ literal پایه در آداپتر عوض نشد. نمونهٔ تحویلی همچنان PB10 است.}

\subsection{مرحلهٔ ۸: توالی سنسور را بدون خواندن کل کد دنبال کنیم}
دو درخواست از سنسور را تصور کنید: «الان اندازه بگیر» و بعد «نتیجه چه بود؟». هر دو با reset شروع می‌شوند تا توجه سنسور جلب شود.
\begin{enumerate}
\item reset و presence بررسی شود. بدون پاسخ، وضعیت نبود قطعه برگردد.
\item فرمان
\lr{0xCC}
یعنی انتخاب تنها سنسور و فرمان
\lr{0x44}
یعنی شروع تبدیل دما ارسال شود.
\item تا پایان تبدیل polling شود و بین busyها انتظار باشد. با تمام‌شدن بودجه، انتظار خاتمه پیدا کند.
\item دوباره reset و فرمان‌های
\lr{0xCC}
و
\lr{0xBE}
برای خواندن scratchpad، یعنی ناحیهٔ ۹ بایتی نتیجه، ارسال شوند.
\item بایت بررسی کنترل و دما decode شود. در موفقیت OK و در خطا، دمای خروجی کاربر بدون تغییر بماند.
\end{enumerate}
فرمان انتخاب عمومی دلیل شرط تک‌سنسور است. چند سنسور هم‌زمان پاسخ می‌دهند و این نسخه تعداد یا آدرس فردی آن‌ها را پیدا نمی‌کند. VDD مستقل، polling حین تبدیل را ممکن می‌کند؛ تغذیهٔ انگلی مدار و روش دیگری می‌خواهد.

دمای خام در دو بایت اول، کم‌ارزش‌تر در ابتدا، قرار دارد. برای بایت‌های
\lr{91 01}
مقدار
\lr{0x0191}
یا ۴۰۱ به دست می‌آید. واحد خام یک‌شانزدهم درجه است؛ ۴۰۱ تقسیم بر ۱۶ برابر ۲۵٫۰۶۲۵ درجه می‌شود. نمونهٔ منفی
\lr{FF5E}
به منفی ۱۶۲ تقسیم بر ۱۶، یعنی منفی ۱۰٫۱۲۵ درجه تبدیل می‌شود.

CRC خرابی بایت‌ها را بررسی می‌کند. بایت‌های ثابت هم کنترل می‌شوند چون ۹ صفر، CRC صفر معتبر دارد؛ فقط CRC خط پایین را نمی‌گیرد. رزولوشن تعداد بیت‌های معنادار پایین را مشخص می‌کند و بیت تعریف‌نشده قبل از تبدیل حذف می‌شود. ۸۵ درجهٔ معتبر پذیرفته می‌شود چون ممکن است دمای واقعی باشد؛ تکرار نامعقول آن، بررسی تغذیه و زمان تبدیل می‌خواهد.

\subsection{مرحلهٔ ۹: تابع قابل‌حمل را بخش‌به‌بخش بخوانیم}
این کل عملیات خواندن دماست؛ فایل C آن هدر MCU ندارد:
\begin{LTR}
\lstinputlisting{../../ds18b20/Libraries/DS18B20/Src/ds18b20.c}
\end{LTR}
بلوک اول اشاره‌گرها و callback را کنترل می‌کند. بعد تبدیل شروع می‌شود. حلقۀ while انتظار محدود دارد. reset دوم خواندن نتیجه را شروع می‌کند. حلقۀ for نه بایت را می‌گیرد. CRC و مقدار ثابت، دادهٔ خراب را رد می‌کنند. در پایان دو بایت ترکیب، علامت منفی و محدوده بررسی و سلسیوس در متغیر کاربر نوشته می‌شود.

۸۰۰ انتظار میلی‌ثانیهٔ اسمی deadline ساعت واقعی نیست. تایم‌اسلات، وقفه و گردشدن تیک HAL سربار دارند؛ timeout نمونه در busy دائمی می‌تواند حدود ۱٫۷ ثانیه شود. فراخوانی blocking است: اجرای اصلی انتظار می‌کشد ولی وقفه بین بخش‌های حساس می‌تواند اجرا شود. در ISR یا با وقفهٔ بسته صدا نزنید.

\subsection{مرحلهٔ ۱۰: وقتی خواستیم منتقل کنیم، آداپتر را بفهمیم}
callback یک تابع عادی است که آدرسش به تابع دیگر داده می‌شود. به‌جای واردکردن رجیستر STM32 به کتابخانه، توابعی با معنی «پایین بکش»، «آزاد کن»، «بخوان» و «صبر کن» به آن می‌دهیم. آداپتر آن‌ها را فراهم می‌کند.
\begin{LTR}
\begin{lstlisting}
/* Inside the STM32 adapter: */
static void low(void *context)
{
    (void)context;
    ONE_WIRE_1_GPIO_Port->BRR = ONE_WIRE_1_Pin;
}
\end{lstlisting}
\end{LTR}
static این تابع کمکی را خصوصی همان فایل C نگه می‌دارد. context در برد دیگر می‌تواند شیء پورت و پایه را منتقل کند؛ نمونهٔ پایهٔ ثابت به آن نیاز ندارد. ماکرو GPIO از CubeMX است، بنابراین تغییر پایه به ویرایش این تابع نیاز ندارد.

هفت تابع، هفت کار جدا دارند: پایین‌کشیدن، آزادکردن، خواندن سطح، تأخیر میکروثانیه، ذخیره و بستن وقفه، بازیابی وقفه و تأخیر میلی‌ثانیه. تنظیم باس آدرس آن‌ها را نگه می‌دارد:
\begin{LTR}
\par\noindent\begin{minipage}{\linewidth}
\lstinputlisting[firstline=43,lastline=60]{../../ds18b20/Application/Src/ds18b20_stm32.c}
\end{minipage}\par
\end{LTR}
شمارندۀ سیکل DWT، تأخیر کوتاه را می‌سازد؛ در ۷۲ مگاهرتز هر میکروثانیه ۷۲ سیکل است. HAL انتظار میلی‌ثانیه را می‌دهد. توابع وقفه وضعیت ذخیره‌شده را برمی‌گردانند و وقفه را بی‌قید فعال نمی‌کنند.

لایهٔ سیم بایت را از کم‌ارزش‌ترین بیت می‌فرستد. reset، ۴۸۰ میکروثانیه پایین می‌کشد و ۷۰ پس از آزادکردن presence را می‌خواند. ارسال یک، ۶ پایین و ۶۴ آزاد؛ صفر، ۶۰ پایین و ۱۰ آزاد است. خواندن، ۳ پایین، ۱۰ پس از آزادکردن تا نمونه و ۵۳ recovery درخواست می‌کند. سربار GPIO و فراخوانی روی برد اندازه‌گیری شود. reset حدود ۵۵۰ میکروثانیه به‌اضافهٔ سربار، وقفه را می‌بندد؛ برای deadline سخت‌تر طراحی زمان دیگری لازم است.

برای MCU دیگر، آداپتر را بنویسید و زمان‌ها را اندازه بگیرید؛ دو فایل C قابل‌حمل حفظ شوند. برای چند سنسور یا تبدیل غیرمسدودکننده ابتدا رفتار API و تست تعریف شود. این قابلیت‌ها عمداً خارج از درس اول هستند.

\subsection{مرحلهٔ ۱۱: برد واقعاً چه چیزی را ثابت کرد؟}
صاحب پروژه گزارش کرده firmware قبلی را با کیل دانلود و اجرا کرده و هنگام breakpoint دما دیده شده است. این بازخورد اولیهٔ سخت‌افزار مفید است، اما هنوز refresh پیوسته، دقت دما، قطع و وصل و زمان موج واقعی را ثابت نمی‌کند. تنظیم اصلاح‌شده باید در جلسۀ واقعی دوباره امتحان شود.

این موارد جدا ثبت شوند:
\begin{enumerate}
\item refresh فعال: موفقیت‌ها بدون breakpoint تراکنش زیاد شوند.
\item گرم و سردکردن: چند دمای موفق متناسب تغییر کنند.
\item قطع و وصل: خطا اعتبار را پاک و اتصال مجدد خواندن را برگرداند.
\item مقایسه با دماسنج مرجع و بررسی DQ زیر بار وقفهٔ واقعی انجام شود.
\item نسخه، ولتاژ، کریستال، مقاومت، سیم، پروگرامر، دما و موج ثبت شوند.
\end{enumerate}
مشاهدات واقعی در فایل وضعیت ثبت شوند:
\lr{PROJECT\_STATUS.md}.
ساخت و شبیه‌سازی شواهد نرم‌افزارند؛ پذیرش برد به مشاهدات بالا نیاز دارد. سه تست C از Python Unicorn و ARM GCC مربوط به PlatformIO استفاده می‌کنند.
